\documentclass[11pt,a4paper]{article}

% --- Encoding ---
\usepackage[utf8]{inputenc}
\usepackage[T1]{fontenc}
\usepackage[spanish,es-noshorthands]{babel}
\usepackage{lmodern}

% --- Layout ---
\usepackage[margin=2.5cm]{geometry}
\usepackage{parskip}
\setlength{\parindent}{0pt}
\setlength{\parskip}{0.5em plus 0.2em minus 0.1em}

% --- Tables and graphics ---
\usepackage{booktabs}
\usepackage{array}
\usepackage{enumitem}
\usepackage{xcolor}
\usepackage{graphicx}
\usepackage{tikz}
\usetikzlibrary{positioning,shapes.geometric,arrows.meta,fit,backgrounds,calc,decorations.pathreplacing}
\usepackage{caption}
\usepackage{amsmath,amssymb}

% --- Code listings ---
\usepackage{listings}
\definecolor{codegray}{rgb}{0.96,0.96,0.96}
\definecolor{codeborder}{rgb}{0.85,0.85,0.85}
\definecolor{codecomment}{rgb}{0.30,0.55,0.30}
\definecolor{codekeyword}{rgb}{0.10,0.30,0.70}
\definecolor{codestring}{rgb}{0.65,0.15,0.10}

\lstdefinestyle{cstyle}{
    language=C,
    backgroundcolor=\color{codegray},
    basicstyle=\ttfamily\footnotesize,
    breaklines=true,
    frame=single,
    rulecolor=\color{codeborder},
    showstringspaces=false,
    keywordstyle=\color{codekeyword}\bfseries,
    commentstyle=\color{codecomment}\itshape,
    stringstyle=\color{codestring},
    aboveskip=1em,
    belowskip=1em,
}

\lstdefinestyle{shellstyle}{
    language=bash,
    backgroundcolor=\color{codegray},
    basicstyle=\ttfamily\footnotesize,
    breaklines=true,
    frame=single,
    rulecolor=\color{codeborder},
    showstringspaces=false,
    keywordstyle=\color{codekeyword}\bfseries,
    commentstyle=\color{codecomment}\itshape,
    stringstyle=\color{codestring},
    aboveskip=1em,
    belowskip=1em,
    columns=fullflexible,
}

\lstdefinestyle{textstyle}{
    backgroundcolor=\color{codegray},
    basicstyle=\ttfamily\footnotesize,
    breaklines=true,
    frame=single,
    rulecolor=\color{codeborder},
    showstringspaces=false,
    aboveskip=1em,
    belowskip=1em,
    columns=fullflexible,
}

\lstset{style=textstyle}

% --- Definition / Idea / Important boxes ---
\usepackage[most]{tcolorbox}
\newtcolorbox{definicion}[1][]{
    colback=blue!5!white,
    colframe=blue!50!black,
    fonttitle=\bfseries,
    title=Definición,
    sharp corners=southwest,
    rounded corners=northwest,
    arc=2pt,
    boxrule=0.6pt,
    left=8pt,right=8pt,top=4pt,bottom=4pt,
    #1
}
\newtcolorbox{idea}[1][]{
    colback=orange!5!white,
    colframe=orange!60!black,
    fonttitle=\bfseries,
    title=Idea intuitiva,
    sharp corners=southwest,
    rounded corners=northwest,
    arc=2pt,
    boxrule=0.6pt,
    left=8pt,right=8pt,top=4pt,bottom=4pt,
    #1
}
\newtcolorbox{importante}[1][]{
    colback=red!5!white,
    colframe=red!50!black,
    fonttitle=\bfseries,
    title=Importante,
    sharp corners=southwest,
    rounded corners=northwest,
    arc=2pt,
    boxrule=0.6pt,
    left=8pt,right=8pt,top=4pt,bottom=4pt,
    #1
}

% --- Hyperlinks ---
\usepackage[colorlinks=true,
            linkcolor=blue!60!black,
            urlcolor=blue!50!black,
            citecolor=blue!60!black]{hyperref}

% --- Color palette ---
\definecolor{fs1}{HTML}{4C72B0}
\definecolor{fs2}{HTML}{DD8452}
\definecolor{fs3}{HTML}{55A868}
\definecolor{fs4}{HTML}{8172B2}
\definecolor{fs5}{HTML}{C44E52}

% --- Title ---
\title{
    \vspace{-2.5em}
    \textbf{Resumen teórico --- Sistemas Operativos} \\[0.3em]
    \large Notas del curso: \emph{Sistemas de archivos}
}
\author{
    Tomás Rodeghiero \\
    \small Sistemas Operativos --- Universidad Nacional de Río Cuarto
}
\date{2026}

\begin{document}

\maketitle

\begin{abstract}
\noindent
Este documento es un resumen teórico del capítulo \emph{Sistemas de
archivos} de las \emph{Notas del curso} de Sistemas Operativos (Marcelo
Arroyo, UNRC). Se presentan la estructura lógica de un disco, la
interfaz de un sistema de archivos (atributos, directorios,
operaciones, llamadas al sistema UNIX), el diseño del \emph{Virtual
File System} (VFS) y el \emph{common file model}, la representación
de archivos y directorios mediante FAT y mediante inodes (incluyendo
inodes indirectos), el mecanismo de montaje, la abstracción de
dispositivos como archivos especiales en \texttt{/dev}, el descriptor
de archivo, las operaciones sobre dispositivos de caracteres, el caso
de estudio del filesystem por capas de xv6 (con su mecanismo
transaccional de logging), el proceso de booting (BIOS, GRUB, UEFI,
secure boot, GPT) y los sistemas de archivos especiales (procfs,
sysfs, ramfs, initramfs, busybox). El objetivo es servir de material
de estudio para preparar los prácticos asociados al sistema de
archivos y el final de la materia.
\end{abstract}

\tableofcontents
\bigskip

% =====================================================================
\section{Introducción y estructura lógica de un disco}
% =====================================================================

La memoria principal o RAM es \textbf{volátil}: su contenido se pierde
al apagar el sistema. Además, tiene capacidad limitada. Por eso se
necesita un medio de almacenamiento \emph{masivo} y \emph{persistente}
para los datos y programas del sistema y de los usuarios.

Los dispositivos de almacenamiento han evolucionado a través del tiempo
(tarjetas/cintas perforadas, discos magnéticos y ópticos, hasta las
actuales memorias no volátiles de estado sólido: pendrives, SSD con
interfaz ATA, etc.). En estas notas se denomina genéricamente
\textbf{disco} a un dispositivo de almacenamiento.

\begin{idea}
Estos dispositivos almacenan datos en \emph{bloques} de tamaño fijo
(por ejemplo, 512\,bytes). El sistema operativo no lee ni escribe
``bytes sueltos'' contra un disco: trabaja siempre con bloques.
\end{idea}

\subsection{Particiones y volúmenes}

Un dispositivo de almacenamiento puede dividirse lógicamente en
\textbf{particiones} o \textbf{volúmenes}. Esto permite implementar
sistemas de archivos que no deban manejar un espacio enorme de números
de bloques. Un dispositivo no particionado provee \emph{al menos un}
volumen.

\subsection{Formato a bajo nivel}

Un disco se formatea a bajo nivel con la siguiente estructura lógica:

\begin{center}
\begin{tikzpicture}[
    every node/.style={font=\ttfamily\small},
    cell/.style={draw,minimum height=0.8cm,align=center},
]
\node[cell,fill=fs5!20,minimum width=0.8cm] (bs) {bs};
\node[cell,fill=fs1!15,minimum width=2.4cm,right=0pt of bs] (pt) {partition table};
\node[cell,fill=fs3!15,minimum width=1.7cm,right=0pt of pt] (v1) {volume 1};
\node[cell,fill=fs2!15,minimum width=1.7cm,right=0pt of v1] (v2) {volume 2};
\node[cell,minimum width=1.5cm,right=0pt of v2] (dots) {\ldots};
\end{tikzpicture}
\end{center}

\begin{itemize}[leftmargin=*]
    \item \textbf{Boot sector (bs).} Contiene un pequeño programa
        encargado de cargar el kernel del SO en memoria y transferirle
        el control.
    \item \textbf{Tabla de partición.} Información (comienzo, fin,
        booteable, tipo de FS, \ldots) sobre la división lógica del
        disco. Cada parte se denomina \emph{partición} o
        \emph{volumen}.
    \item \textbf{Volúmenes.} En cada uno se almacena un sistema de
        archivos (FS).
\end{itemize}

\begin{importante}
El estándar BIOS incluye la tabla de particiones primaria (de
\emph{cuatro entradas}) dentro del boot sector. El estándar moderno
UEFI se basa en una estructura mucho menos restrictiva, descripta más
adelante.
\end{importante}

Una plataforma de hardware puede incluir varios dispositivos de
almacenamiento de distintas tecnologías. El subsistema del SO que
provee abstracciones para gestionar objetos de datos y programas se
conoce como \textbf{sistema de archivos (filesystem)}.

% =====================================================================
\section{Interfaz de un sistema de archivos}
% =====================================================================

El sistema de archivos abstrae una unidad de almacenamiento en un
\textbf{archivo}. Para organizar la enorme cantidad de archivos de un
sistema, éstos se almacenan en contenedores llamados
\textbf{directorios} o carpetas.

\begin{definicion}
Un \textbf{archivo} se representa como un par
$(\textit{attributes},\;\textit{data})$: una colección de
\emph{metadatos} (atributos) más una secuencia lógicamente contigua de
\emph{datos}, aunque físicamente esos datos pueden estar dispersos en
distintos bloques del dispositivo.
\end{definicion}

\subsection{Atributos típicos}

\begin{itemize}[leftmargin=*,nosep]
    \item \textbf{Nombre.}
    \item \textbf{Tipo} (según el formato lógico del contenido).
    \item \textbf{Tamaño} (comúnmente en bytes).
    \item \textbf{Dueño}: usuario que define el control de acceso.
    \item \textbf{Creador}: usuario que lo creó originalmente.
    \item \textbf{Permisos de acceso}: qué usuarios pueden acceder y
        cómo (lectura, escritura, ejecución, etc.).
    \item \textbf{Hora y fecha de último acceso}; otros tiempos
        (creación, modificación).
\end{itemize}

\subsection{Árbol de directorios y paths}

Los directorios contienen archivos y \emph{subdirectorios}, formando
una estructura de \textbf{árbol}: los archivos son las hojas y los
directorios los nodos intermedios. Todo sistema de archivos contiene
al menos un \textbf{directorio raíz}:

\begin{itemize}[leftmargin=*,nosep]
    \item En SO tipo UNIX se denota \texttt{/}.
    \item En MS-Windows se denota \texttt{\textbackslash}
        (\emph{backslash}).
\end{itemize}

Para identificar unívocamente a un archivo se usa su \textbf{path}
(camino desde la raíz), por ejemplo
\verb|/home/os-student/books/SO-book.pdf|. Esto permite reusar nombres
en distintas carpetas.

\subsection{Tipos de archivo}

En DOS o MS-Windows, el tipo del archivo se determina por su
\emph{extensión} (sufijo del nombre): \verb|.com|, \verb|.exe|, etc.
En los SO tipo UNIX, el tipo está en los \emph{metadatos} y la
extensión es sólo una convención entre usuarios; el sistema no la
interpreta (por ejemplo, \verb|ld-linux.so.2|). UNIX clasifica los
archivos en:

\begin{itemize}[leftmargin=*]
    \item \textbf{Regulares}: archivos de datos y ejecutables.
    \item \textbf{Directorios}: contenedores de archivos y
        subdirectorios.
    \item \textbf{Enlaces simbólicos}: referencias (alias) a otros
        archivos.
    \item \textbf{Pipes y FIFOs}: usados para IPC.
    \item \textbf{Especiales}: representan dispositivos (de bloques o
        caracteres).
    \item \textbf{Sockets}: mecanismo de comunicación local o remoto.
\end{itemize}

% =====================================================================
\section{Operaciones sobre archivos y directorios}
% =====================================================================

\subsection{Operaciones sobre archivos regulares}

\begin{itemize}[leftmargin=*,nosep]
    \item Creación, lectura, escritura, ejecución, borrado
        (\emph{delete}).
    \item \textbf{Append}: escribir datos al final.
    \item \textbf{Truncado}: borrar datos al final.
    \item \textbf{Posicionado}: mover el cursor de lectura/escritura.
    \item \textbf{Renombrado}: cambiar el nombre.
    \item \textbf{Locking}: obtener acceso exclusivo.
    \item Cambiar otros metadatos (dueño, permisos, \ldots).
\end{itemize}

\subsection{Operaciones sobre directorios}

\begin{itemize}[leftmargin=*]
    \item \textbf{Creación.} Es común que un directorio vacío se cree
        con \emph{dos enlaces}: \verb|..| (referencia al directorio
        padre) y \verb|.| (referencia al directorio actual).
    \item \textbf{Búsqueda} de un archivo o subdirectorio. Cada
        elemento dentro de un directorio se llama \emph{directory
        entry}.
    \item \textbf{Eliminación.} Comúnmente se exige que el directorio
        esté vacío.
\end{itemize}

\subsection{Paths absolutos y relativos}

En las operaciones que reciben un \texttt{path}, éste puede ser un
\textbf{full path} (desde la raíz) o \textbf{relativo} al directorio
corriente. Por ejemplo:

\begin{itemize}[leftmargin=*,nosep]
    \item \verb|open("/home/user1/docs/mycv.tex")| usa un full path.
    \item \verb|open("docs/mycv.tex")| refiere al mismo archivo si el
        directorio corriente es \verb|/home/user1/|.
\end{itemize}

\paragraph{Variables de ambiente.}
Un programa se ejecuta en un \emph{ambiente} representado por un
conjunto de \textbf{variables de ambiente}, cada una un par
\verb|variable=value|. En UNIX, un programa accede a ellas mediante el
tercer argumento de \texttt{main}. El directorio corriente está
representado por la variable \texttt{PWD} (\verb|echo $PWD|).

\subsection{Llamadas al sistema (UNIX)}

Algunas llamadas al sistema sobre archivos en UNIX:

\begin{lstlisting}[style=cstyle]
int  creat (file_path, flags);             // creacion
int  open  (file_path, flags);             // apertura -> file descriptor
int  close (fd);                           // cierre
int  read  (fd, buffer, count);            // lectura desde el cursor
int  write (fd, buffer, count);            // escritura desde el cursor
int  seek  (fd, pos, from);                // posicionado del cursor
int  unlink(path);                         // eliminacion
int  link  (dst, name);                    // hard link
int  symlink(dst, name);                   // symbolic link
int  stat  (fd, stat_buffer);              // metadatos
\end{lstlisting}

Y sobre directorios:

\begin{lstlisting}[style=cstyle]
int mkdir(path, mode);
int rmdir(path);
\end{lstlisting}

\begin{idea}
\texttt{open} retorna un \emph{file descriptor} que se usa en todas las
operaciones siguientes. Posiciona el cursor en cero al abrir en modo
lectura/escritura, o al final si se abre en modo \texttt{APPEND}.
\texttt{read} retorna 0 cuando el cursor estaba en EOF.
\end{idea}

Para acceder al \emph{contenido} de un directorio, éste puede leerse
como un archivo cuyos registros son los metadatos de los archivos que
contiene. La biblioteca estándar de cada SO ofrece funciones para
manipular directorios (en UNIX, ver \texttt{readdir()}).

% =====================================================================
\section{Virtual File System (VFS)}
% =====================================================================

Los SO modernos soportan diferentes tipos (o formatos) de sistemas de
archivos en distintos dispositivos y tecnologías. En estos casos debe
presentarse al usuario un sistema de archivos \emph{uniforme} con
operaciones que oculten los detalles de bajo nivel.

\begin{definicion}
El \textbf{Virtual File System (VFS)} es el subsistema del kernel
que, ante una llamada al sistema sobre archivos, determina
\emph{sobre qué dispositivo} y \emph{con qué sistema de archivos
concreto} debe realizarse la operación. Define un \emph{common file
model} basado --- en los UNIX --- en nodos índices (\emph{inodes}).
\end{definicion}

\subsection{Diseño del VFS}

\begin{center}
\begin{tikzpicture}[
    every node/.style={font=\ttfamily\small},
    box/.style={draw,minimum height=0.7cm,align=center},
    big/.style={draw,minimum height=0.9cm,align=center}
]
\node[big,fill=fs1!20,minimum width=1.8cm] (proc) at (0,5.0) {process};
\node[anchor=west] at (proc.east) {syscalls};

\node[big,fill=fs3!15,minimum width=8.5cm,minimum height=1.3cm] (vfs) at (0,3.4) {};
\node[font=\bfseries\small] at (0,3.85) {VFS - common file model};
\node[box,fill=white,minimum width=0.9cm] (sb)    at (-3.2,3.05) {sb};
\node[box,fill=white,minimum width=0.9cm] (file)  at (-1.9,3.05) {file};
\node[box,fill=white,minimum width=1.2cm] (inode) at (-0.4,3.05) {inode};
\node[box,fill=white,minimum width=0.9cm] (dir)   at (1.2,3.05) {dir};
\node[box,fill=white,minimum width=1.0cm] (fops)  at (2.6,3.05) {fops};

\node[box,fill=fs2!15] (ext3) at (-3.5,1.7) {ext3};
\node[box,fill=fs2!15] (fat32) at (-1.75,1.7) {fat32};
\node[box,fill=fs2!15] (procfs) at (0,1.7) {proc};
\node[box,fill=fs2!15] (ntfs) at (1.75,1.7) {NTFS};
\node[box,fill=fs2!15] (nfs) at (3.5,1.7) {NFS};

\node[big,fill=fs4!20,minimum width=3.0cm] (bc) at (0,0.4) {buffer cache};
\node[big,fill=fs5!20,minimum width=4.5cm] (bdd) at (0,-0.9) {block device drivers};

\draw[thick] (proc.south) -- (vfs.north);
\draw[thick] (ext3.north)   -- ($(ext3.north)+(0,0.5)$);
\draw[thick] (fat32.north)  -- ($(fat32.north)+(0,0.5)$);
\draw[thick] (procfs.north) -- ($(procfs.north)+(0,0.5)$);
\draw[thick] (ntfs.north)   -- ($(ntfs.north)+(0,0.5)$);
\draw[thick] (nfs.north)    -- ($(nfs.north)+(0,0.5)$);

\draw[thick] (ext3.south) -- (bc.north);
\draw[thick] (fat32.south) -- (bc.north);
\draw[thick] (procfs.south) -- (bc.north);
\draw[thick] (ntfs.south) -- (bc.north);
\draw[thick] (nfs.south) -- (bc.north);

\draw[thick] (bc.south) -- (bdd.north);
\end{tikzpicture}
\end{center}

\subsection{Estructuras de datos del common file model}

El VFS define una interfaz que cada FS concreto debe implementar y
estructuras de datos comunes que ocultan los detalles de cada filesystem
real. Cada implementación concreta debe \textbf{mapear} su representación
al modelo abstracto.

\begin{itemize}[leftmargin=*]
    \item \textbf{Superblock (sb).} Metadatos del filesystem: tipo o
        id, número y tamaño de bloques, ubicación del directorio raíz,
        etc.
    \item \textbf{File.} Representa un archivo \emph{abierto por un
        proceso}. Contiene al menos un puntero al \emph{inode} del
        archivo y un \emph{offset} (cursor de lectura/escritura).
    \item \textbf{Inode.} Contiene los \emph{metadatos} del archivo:
        tamaño, dispositivo, bloques de datos usados, etc.
    \item \textbf{Directory (dir).} Un directorio es un archivo de
        datos cuyas entradas son pares (\emph{name}, \emph{inode
        number}). El VFS define también una interfaz abstracta para
        buscar por \emph{file paths}.
    \item \textbf{fops.} Interfaz abstracta que cada filesystem
        concreto debe implementar con operaciones como
        \texttt{open}, \texttt{read}/\texttt{write} y
        \texttt{free} sobre inodes.
\end{itemize}

\subsection{Servicios comunes provistos por el VFS}

\begin{itemize}[leftmargin=*]
    \item Punto de entrada de las llamadas al sistema (\texttt{open},
        \texttt{read}, \texttt{write}, \ldots).
    \item \textbf{Buffers de caché} para dispositivos de bloques. Los
        bloques de datos se mantienen en una caché con \emph{pool}
        limitado y políticas de reemplazo (LRU, similar a memoria
        virtual). \emph{No} se usa caché sobre dispositivos de
        caracteres.
    \item \textbf{Gestión de los requerimientos de E/S}, basada en
        colas y algoritmos de ordenamiento (algoritmo del elevador
        para discos, FIFO para ttys).
    \item Registración de los sistemas de archivos concretos.
    \item Registración de los device drivers.
    \item Montaje y desmontaje de filesystems concretos.
\end{itemize}

\begin{importante}
El módulo de \emph{buffers de caché} es el que se comunica con los
\emph{block device drivers}. Los filesystems usan sus servicios para
leer/escribir bloques: nunca tocan directamente al driver.
\end{importante}

% =====================================================================
\section{Representación de archivos y directorios}
% =====================================================================

Para gestionar archivos y directorios se necesitan estructuras de
control y metadatos almacenadas en el propio dispositivo. Mínimamente:

\begin{itemize}[leftmargin=*,nosep]
    \item Metadatos del FS (ubicación del bloque del directorio raíz,
        número de bloques en el dispositivo, etc.).
    \item Conjunto de bloques libres y usados.
    \item Descripción de espacios de bloques de datos y de metadatos.
\end{itemize}

La forma concreta que toman esas estructuras se llama \emph{formato del
sistema de archivos}. Crear un sistema de archivos inicial se conoce
como \textbf{formatear} el dispositivo (comandos \texttt{mkfs} en UNIX
o \texttt{format} en DOS/Windows).

A continuación se analizan dos formatos ampliamente usados: FAT e
inodes.

% ---------------------------------------------------------------------
\subsection{FAT}
% ---------------------------------------------------------------------

\textbf{FAT} (\emph{File Allocation Table}) fue desarrollado para DOS y
se usó en las primeras versiones de MS-Windows. Su estructura:

\begin{center}
\begin{tikzpicture}[
    every node/.style={font=\ttfamily\small},
    cell/.style={draw,minimum height=0.7cm,align=center},
]
\node[cell,fill=fs5!15,minimum width=0.7cm] (sb) {sb};
\node[cell,fill=fs1!15,minimum width=1.2cm,right=0pt of sb] (fat) {FAT};
\node[cell,fill=fs3!15,minimum width=1.2cm,right=0pt of fat] (root) {root};
\node[cell,minimum width=1.0cm,right=0pt of root] (blk1) {};
\node[cell,minimum width=2.6cm,right=0pt of blk1] (dots) {\ldots};
\node[cell,minimum width=0.8cm,right=0pt of dots] (blkN) {};

\draw[decorate,decoration={brace,mirror,amplitude=4pt}]
    ($(blk1.south west)+(0,-0.1)$) -- ($(blkN.south east)+(0,-0.1)$)
    node[midway,below=4pt,font=\small] {data blocks (clusters)};
\end{tikzpicture}
\end{center}

\paragraph{Superblock.}
Contiene metadatos como el tipo de FS (\emph{FAT magic number}), el
número de bloques del volumen, el tamaño de bloque, etc.

\paragraph{Clusters.}
Para reducir el número de bloques y la fragmentación de datos, los
bloques físicos se agrupan en \textbf{clusters}: grupos de bloques
contiguos (típicamente de 1, 2, 4, 8 o 16 bloques). Un cluster es la
\emph{unidad mínima} de asignación de datos para un archivo. El primer
cluster de datos contiene el directorio raíz. En lo que sigue, ``bloque''
y ``cluster'' se usan como sinónimos.

\paragraph{La tabla FAT.}
\texttt{fat[N]} es un arreglo con tantas entradas como clusters de
datos. Implementa \emph{listas enlazadas} de los bloques de datos de
cada archivo. Un directorio se representa como un archivo cuya
secuencia de registros tiene la forma:

\begin{center}
\begin{tikzpicture}[
    every node/.style={font=\ttfamily\small},
    cell/.style={draw,minimum height=0.7cm,align=center},
]
\node[cell,minimum width=2.4cm,fill=fs1!15] (name) {file name};
\node[cell,minimum width=1.2cm,fill=fs2!15,right=0pt of name] (a) {attrs};
\node[cell,minimum width=1.6cm,fill=fs3!15,right=0pt of a] (c) {cluster1};
\end{tikzpicture}
\end{center}

En FAT32 cada entrada tiene 32\,bits con 12\,bytes para el nombre y
bits de atributos (\emph{size}, \emph{read only}, \emph{hidden},
fechas, \ldots).

\paragraph{Lectura de un archivo.}
\texttt{cluster1} contiene el número del primer cluster del archivo.
Los siguientes clusters se obtienen indexando en la FAT:
\verb|FAT[cluster1]| contiene el \emph{next cluster}. Un valor 0
indica fin del archivo. Así, un archivo con tres clusters
$j \to i \to ?$ se representa con
$\textit{FAT}[j] = i$, $\textit{FAT}[i] = 0$ (terminador).

\begin{idea}
La FAT contiene la lista de clusters de cada archivo. Los clusters
\emph{libres} pueden inferirse de la FAT: si un cluster no es
alcanzable desde ningún directorio, está libre. No hace falta
almacenarlos en disco; pueden calcularse al inicio y mantenerse en
memoria.
\end{idea}

% ---------------------------------------------------------------------
\subsection{Nodos índices (inodes)}
% ---------------------------------------------------------------------

Una técnica alternativa, usada por la mayoría de los FS para grandes
volúmenes (\textbf{ext2/3/4} en GNU/Linux y otros UNIX), \emph{separa}
los metadatos de la entrada del directorio. Un \textbf{inode}
(\emph{index node}) contiene los metadatos del archivo más la lista de
bloques de datos que usa.

\paragraph{Estructura del filesystem.}

\begin{center}
\begin{tikzpicture}[
    every node/.style={font=\ttfamily\small},
    cell/.style={draw,minimum height=0.7cm,align=center},
]
\node[cell,fill=fs5!15,minimum width=0.7cm] (sb) {sb};
\node[cell,fill=fs1!15,minimum width=1.2cm,right=0pt of sb] (bm) {bitmap};
\node[cell,fill=fs4!15,minimum width=1.4cm,right=0pt of bm] (in) {inodes};
\node[cell,minimum width=1.0cm,right=0pt of in] (b1) {};
\node[cell,minimum width=2.4cm,right=0pt of b1] (dots) {\ldots};
\node[cell,minimum width=0.7cm,right=0pt of dots] (bN) {};

\draw[decorate,decoration={brace,mirror,amplitude=4pt}]
    ($(b1.south west)+(0,-0.1)$) -- ($(bN.south east)+(0,-0.1)$)
    node[midway,below=4pt,font=\small] {data blocks (clusters)};
\end{tikzpicture}
\end{center}

\paragraph{Contenido de un inode.}
Un inode contiene:

\begin{itemize}[leftmargin=*,nosep]
    \item Atributos (dueño, permisos, tamaño en bytes, fechas, \ldots).
    \item Un arreglo o tabla de números de bloques de datos
        (\emph{data blocks table}) de tamaño \emph{predefinido} $k$,
        que representa la secuencia de bloques que ocupa el archivo.
\end{itemize}

\paragraph{Directorios.}
Un directorio es un archivo estructurado cuyas entradas tienen sólo
dos campos: $(\textit{filename},\;\textit{inode})$. Habitualmente el
directorio raíz está representado por el primer inode. Los inodes
están numerados/indexados.

\paragraph{Gestión de espacios libres.}
Es conveniente usar un \textbf{bitmap} para la gestión de inodes y
bloques de datos libres. Si bien se podrían deducir por
\emph{complemento} de los alcanzables, esto sería costoso porque la
información está distribuida en todos los inodes.

\paragraph{Tamaño de un inode.}
Por eficiencia se elige el tamaño de un inode igual al de un bloque
físico (sector) del disco.

\begin{importante}
Un inode con tabla de tamaño $k$ puede representar archivos de \emph{a
lo sumo} $k$ bloques de datos. Para archivos más grandes se usa la
técnica de los \textbf{inodes indirectos}.
\end{importante}

\paragraph{Inodes indirectos.}
La idea es que las últimas entradas de la tabla del inode apunten a
\emph{otros inodes} que sólo contienen tablas de números de bloques
adicionales. Comúnmente las entradas finales son:

\begin{itemize}[leftmargin=*,nosep]
    \item \textbf{Indirecto simple}: árbol de un nivel.
    \item \textbf{Indirecto doble}: árbol de dos niveles.
    \item \textbf{Indirecto triple}: árbol de tres niveles.
\end{itemize}

Cuantos más niveles de indirección, más grande puede ser el archivo,
pero más lentos son los accesos a los últimos bloques (más
indirecciones a resolver).

\paragraph{En la creación del filesystem.}
Al ejecutar \texttt{mkfs} se decide la cantidad de inodes a reservar en
el disco. El número de inodes determina la \emph{cantidad máxima de
archivos y directorios} que puede tener el filesystem.

% =====================================================================
\section{Montaje de sistemas de archivos}
% =====================================================================

Los SO presentan el espacio de nombres global de dos maneras distintas:

\begin{itemize}[leftmargin=*]
    \item \textbf{MS-Windows.} Cada FS aparece de forma
        \emph{independiente}: un full path toma la forma de un nombre
        de dispositivo (una letra) seguido por dos puntos y un path.
        Típicamente \texttt{C} es el disco primario.
    \item \textbf{UNIX.} Se presenta al usuario un \textbf{único árbol}
        de archivos. La raíz corresponde al FS del dispositivo de
        almacenamiento primario. Para acceder a otro dispositivo hay
        que \textbf{montarlo} en algún directorio (por ejemplo
        \texttt{/mnt/mypendrive}) del árbol. Los comandos son
        \texttt{mount} y \texttt{umount}.
\end{itemize}

\begin{idea}
El mecanismo de montaje permite que las aplicaciones accedan a
archivos con nombres preestablecidos sin nombrar dispositivos
(\texttt{C:}, \texttt{D:}). Esto facilita escribir aplicaciones
portables.
\end{idea}

\subsection{Ejemplo: montar un pendrive FAT en Linux}

\begin{lstlisting}[style=shellstyle]
mkdir /mnt/mypendrive
mount /dev/sdb /mnt/mypendrive
\end{lstlisting}

Esta operación ``cuelga'' el FS del pendrive en el directorio
\verb|/mnt/mypendrive| del FS raíz. Luego, los archivos del pendrive
son accesibles mediante el path \verb|/mnt/mypendrive/...|.

\subsection{Tabla de montaje}

El VFS mantiene una tabla de filesystems montados, por ejemplo:

\begin{center}
\small
\begin{tabular}{@{}llllc@{}}
\toprule
\textbf{mount point} & \textbf{fs type} & \textbf{device} & \textbf{major} & \textbf{access mode} \\
\midrule
\texttt{/}                & ext4    & \texttt{/dev/sda1} & 1 & rw \\
\texttt{/mnt/mypendrive}  & fat32   & \texttt{/dev/sdb}  & 5 & r \\
\texttt{/proc}            & procfs  & proc               & 0 & rw \\
\ldots                    & \ldots  & \ldots             & \ldots & \ldots \\
\bottomrule
\end{tabular}
\end{center}

En UNIX, esta información se obtiene con el comando \texttt{mount} o
leyendo \verb|/etc/mtab|.

\paragraph{umount.}
El comando \verb|umount /mnt/mypendrive| desmonta el filesystem. Es
seguro retirar el dispositivo recién \emph{después} de desmontarlo
porque hace un \texttt{flush} de los bloques que aún pudieran estar en
caché.

\subsection{Interfaz de FS que el VFS exige}

Cada FS concreto debe implementar y registrar:

\begin{itemize}[leftmargin=*,nosep]
    \item \texttt{mount}: cargar y registrar el superblock.
    \item \texttt{umount}: liberar el superblock.
    \item Operaciones sobre directorios (creación, destrucción,
        búsqueda por file paths, \ldots).
    \item Operaciones sobre archivos (inodes).
\end{itemize}

Cada proceso registra sus archivos abiertos en una tabla de su propio
descriptor (en xv6 es el campo \verb|ofile[]| de \verb|struct proc|),
representada como un arreglo de punteros a entradas en la tabla de
\verb|struct file|.

\subsection{Anatomía de un \texttt{open}}

Cuando un proceso ejecuta
\verb|fd = open("/home/user1/docs/mycv.tex", ...)|, el VFS:

\begin{enumerate}[leftmargin=*]
    \item Determina en qué FS concreto está el archivo (en este caso
        el FS raíz, de tipo \texttt{ext4}).
    \item Invoca una función del tipo
        \verb|inode = fs.dir_lookup("/home/user1/docs/mycv.tex")| que
        busca y lee el inode correspondiente.
    \item Crea una nueva entrada en la tabla de archivos abiertos del
        sistema.
    \item Registra el archivo abierto en el descriptor del proceso.
    \item Retorna el índice en \verb|ofiles[]| (el \emph{file
        descriptor}).
\end{enumerate}

\subsection{Anatomía de un \texttt{write}}

Cuando luego el proceso ejecuta \verb|write(fd, buffer, count)|, el
VFS:

\begin{enumerate}[leftmargin=*]
    \item Accede a la \texttt{struct file} correspondiente a
        \texttt{fd}.
    \item De la tabla de montaje determina el FS concreto a usar (por
        ejemplo, el módulo \texttt{ext4}).
    \item Invoca una función del tipo
        \verb|fs.write(inode, buffer, count, offset)|.
        Esta requerirá escribir en uno o más bloques, servicio que
        provee la capa de \emph{buffer de caché} mediante
        \verb|bio_write(buf, dev)|.
    \item Si \texttt{buf} no está en la caché, hay que cargarlo del
        disco en un nuevo buffer; luego se modifica y se marca como
        \emph{dirty} (se escribirá a disco en un \texttt{close(fd)} o
        cuando expire un timer).
    \item Si hay que hacer un requerimiento a disco, éste se encola
        según una política de scheduling (algoritmo del elevador).
    \item El \emph{I/O scheduler} eventualmente invoca a
        \verb|block_dev.do_request(request)| del device driver
        registrado. El \emph{major number} del inode determina cuál es
        el driver; el \emph{minor number} le dice al driver con qué
        dispositivo concreto interactuar.
\end{enumerate}

\begin{idea}
Como un device driver puede controlar varios dispositivos del mismo
tipo, la pareja (major, minor) es la clave: major elige el driver y
minor elige el dispositivo concreto dentro de ese driver.
\end{idea}

% =====================================================================
\section{Dispositivos como archivos especiales}
% =====================================================================

En muchos SO los dispositivos aparecen en el sistema de archivos como
\textbf{archivos especiales}. En UNIX viven bajo \verb|/dev| (por
\emph{devices}). Estos archivos le permiten al VFS mantener la relación
entre los dispositivos y los device drivers correspondientes.

\paragraph{Major y minor numbers.}
Cada archivo especial contiene un par de campos numéricos:

\begin{itemize}[leftmargin=*,nosep]
    \item \textbf{major number}: determina el \emph{tipo} de
        dispositivo. Se usa para relacionar una clase de dispositivos
        con su device driver.
    \item \textbf{minor number}: identifica un dispositivo dentro de
        ese tipo.
\end{itemize}

\subsection{Tabla de device drivers registrados}

Cada device driver se registra en el VFS en su inicio, indicando el
archivo especial y el major number. Así, el VFS mantiene una tabla:

\begin{center}
\begin{tikzpicture}[
    every node/.style={font=\ttfamily\small},
    cell/.style={draw,minimum height=0.55cm,align=center},
]
\node[font=\bfseries\small] at (0,2.5) {devices};
\node[cell,minimum width=1.8cm,fill=fs1!10] (d1) at (0,2.0) {\ldots};
\node[cell,minimum width=1.8cm,fill=fs2!20] (di) at (0,1.45) {};
\node[cell,minimum width=1.8cm,fill=fs1!10] (d3) at (0,0.9) {\ldots};
\node[anchor=east] at (di.west) {i};

\node[font=\bfseries\small] at (4.0,2.5) {dev fops};
\node[cell,minimum width=1.4cm,fill=fs3!20] (op1) at (4.0,2.0) {open};
\node[cell,minimum width=1.4cm,fill=fs3!20,below=0pt of op1] (op2) {write};
\node[cell,minimum width=1.4cm,fill=fs3!20,below=0pt of op2] (op3) {read};
\node[cell,minimum width=1.4cm,below=0pt of op3] (op4) {\ldots};

\node[anchor=west] at (6,2.0) {driver open function};
\node[anchor=west] at (6,1.45) {driver write function};
\node[anchor=west] at (6,0.9) {driver read function};

\draw[-Latex,thick] (di.east) -- ($(op2.west)+(-0.2,0)$);
\draw[-Latex,thick] (op1.east) -- (6,2.0);
\draw[-Latex,thick] (op2.east) -- (6,1.45);
\draw[-Latex,thick] (op3.east) -- (6,0.9);
\end{tikzpicture}
\end{center}

Cada entrada se indexa por el major number. Asocia el dispositivo con
las operaciones registradas por el driver.

% =====================================================================
\section{Descriptor de archivo}
% =====================================================================

\texttt{open} retorna un \textbf{file descriptor} que queda registrado
en la \emph{opened files table} (\verb|ofiles|) del descriptor del
proceso. Esta tabla es un arreglo de punteros a estructuras como:

\begin{lstlisting}[style=cstyle]
struct file {
    union {
        struct pipe  *pipe;
        struct inode *inode;
    };
    int          ref_count;   // procesos trabajando con este file
    unsigned int offset;      // posicion actual de lectura/escritura
};
\end{lstlisting}

Un inode puede representarse como:

\begin{lstlisting}[style=cstyle]
struct inode {
    unsigned int num;        // numero de inode (en disco)
    int          ref_count;
    // los campos siguientes se traen/guardan de/a disco
    short int    type;       // file, directory, FIFO, device, ...
    short int    major;
    short int    minor;
    short int    links;      // numero de symlinks a este inode
    // otros metadatos: owner id, group id, permisos, ...
    unsigned int size;       // en bytes
    unsigned int data_blocks[N];
};
\end{lstlisting}

% =====================================================================
\section{Operaciones con dispositivos de caracteres}
% =====================================================================

La E/S con dispositivos de caracteres se diferencia de la de bloques en
dos aspectos:

\begin{enumerate}[leftmargin=*,nosep]
    \item No tienen filesystem asociado.
    \item La E/S \textbf{no pasa por la caché de bloques} (no es
        buffered): se realiza \emph{directamente} contra el
        dispositivo.
\end{enumerate}

\subsection{Ejemplo: escribir a stdout}

Cuando un proceso desea escribir en el descriptor 1 (salida estándar),
ese descriptor refiere a una \texttt{struct file} cuyo \texttt{inode}
describe un \emph{archivo especial} (por ejemplo \verb|/dev/tty| en
Linux o \texttt{console} en xv6) representando un dispositivo de
caracteres. El \emph{major number} identifica al driver con el que se
debe interactuar. El VFS invoca a la función correspondiente:

\begin{lstlisting}[style=cstyle]
devices[p->ofiles[1]->inode.major].write(p->ofiles[1], ...);
\end{lstlisting}

\subsection{Caso xv6}

En xv6 (similar a UNIX) los character device drivers se registran en
\verb|struct devsw devsw[]|. En xv6-riscv hay un único dispositivo
registrado: la \textbf{consola}. La comunicación se hace por el puerto
serie UART:

\begin{itemize}[leftmargin=*,nosep]
    \item \texttt{uart.c} implementa el \emph{bottom half} (bajo
        nivel).
    \item \texttt{console.c} implementa la lógica de alto nivel
        (\emph{tty discipline}) del driver.
\end{itemize}

El archivo \texttt{console} en la raíz del FS es un archivo especial
que representa al dispositivo. Es abierto por \texttt{init}
\textbf{tres veces} para entrada estándar, salida estándar y salida
estándar de errores. Esos descriptores son \emph{heredados} por los
demás procesos (shell y otros) vía \texttt{fork()}.

% =====================================================================
\section{Filesystem de xv6}
% =====================================================================

Xv6 monta un \emph{único} sistema de archivos con su propio formato
basado en inodes. Se construye al compilar los procesos de usuario
mediante la utilidad \texttt{mkfs} incluida en el build-system.

\subsection{Implementación por capas}

\begin{center}
\begin{tikzpicture}[
    every node/.style={font=\ttfamily\small},
    cell/.style={draw,minimum height=0.6cm,minimum width=3.5cm,align=center},
]
\node[cell,fill=fs1!20] (fd) at (0,4.2) {File descriptor};
\node[cell,fill=fs2!20,below=0pt of fd] (pn) {Pathname};
\node[cell,fill=fs2!20,below=0pt of pn] (di) {Directory};
\node[cell,fill=fs2!20,below=0pt of di] (in) {Inode};
\node[cell,fill=fs3!20,below=0pt of in] (lg) {Logging};
\node[cell,fill=fs4!20,below=0pt of lg] (bc) {Buffer cache};
\node[cell,fill=fs5!20,below=0pt of bc] (dd) {Disk driver};

\node[anchor=west,xshift=4pt] at (fd.east) {file.h/c};
\node[anchor=west,xshift=4pt] at (di.east) {fs.h/c};
\node[anchor=west,xshift=4pt] at (lg.east) {log.c};
\node[anchor=west,xshift=4pt] at (bc.east) {bio.c};
\node[anchor=west,xshift=4pt] at (dd.east) {virtio disc.c};
\end{tikzpicture}
\end{center}

\subsection{FS transaccional con logging}

\begin{importante}
El filesystem de xv6 es \textbf{transaccional} para evitar
inconsistencias en \emph{crashes} o cortes de energía. Antes de escribir
un conjunto de requerimientos a disco, xv6 los escribe en un
\textbf{log} (área especial). Luego realiza un \emph{commit},
replicando (o moviendo) las operaciones del log a los bloques
correspondientes. Finalmente borra el log.
\end{importante}

\paragraph{Recuperación tras crash.}
Cuando el sistema inicia, el FS verifica el log:

\begin{itemize}[leftmargin=*,nosep]
    \item Si está marcado como \emph{completado}, se realizan las
        copias a los bloques de datos contenidos en el log.
    \item Si no, el log se ignora y se elimina. En este caso el
        sistema pierde las últimas escrituras, pero queda en un
        estado \textbf{consistente}.
\end{itemize}

\subsection{Buffer cache simplificada}

El buffer de caché de xv6 consiste en una \textbf{lista doblemente
encadenada} de bloques (buffers). Los buffers son fijos (arreglo
declarado estáticamente) y \emph{no} implementa una técnica de
reemplazo. El manual de xv6 describe en mayor detalle el diseño y la
implementación.

% =====================================================================
\section{Booting}
% =====================================================================

El proceso de inicio (\emph{boot}) comúnmente comienza por un programa
en ROM conocido como \textbf{firmware}.

\subsection{BIOS}

En la PC el firmware se denomina \textbf{Basic Input-Output System
(BIOS)}. Consiste en un conjunto de programas que verifica la RAM
física instalada y algunos dispositivos básicos (teclado, discos).
Suele almacenarse en ROM o EPROM. Usa una memoria volátil mantenida
con una pequeña batería, conocida como \textbf{CMOS}, donde guarda
fecha, hora y parámetros de configuración del hardware.

\paragraph{Carga del kernel.}
La BIOS carga en memoria el \emph{primer sector} del disco primario
(el \textbf{boot sector}) y le transfiere el control. Este programa
\textbf{boot loader} se encarga de cargar en memoria el código y los
datos del kernel del SO y transferirle el control.

El boot loader es muy pequeño (debe caber en un bloque de típicamente
512\,bytes); ese bloque también puede contener la tabla de particiones.
El programa debe saber dónde se encuentra el kernel en el disco o
volumen.

\subsection{Multiboot loaders (GRUB)}

Un \textbf{boot loader genérico multiboot} como \textbf{GRUB} permite
iniciar diferentes sistemas operativos instalados en distintos
dispositivos o volúmenes. Realiza la tarea en \emph{dos etapas}:

\begin{enumerate}[leftmargin=*]
    \item El boot loader del boot sector carga el multiboot loader
        (GRUB) y le transfiere el control.
    \item GRUB --- que conoce formatos de filesystems y de
        ejecutables (ELF, COFF, etc.) --- lee un archivo de
        configuración y presenta un menú al usuario para que
        seleccione el SO a iniciar.
\end{enumerate}

Estos multiboot loaders se usan para iniciar SO tipo UNIX (como
GNU/Linux), donde el kernel es un archivo más cuya ubicación dentro del
FS no se conoce a priori.

El método descripto se conoce como \textbf{BIOS booting} y se
desarrolló en 1975, principalmente para la PC y DOS.

\subsection{UEFI}

Actualmente es común el estándar \textbf{Unified Extensible Firmware
Interface (UEFI)}, basado en un mecanismo mucho más flexible. Sus
características:

\begin{itemize}[leftmargin=*]
    \item Para que sea booteable, un disco debe contener una partición
        de tipo FAT rotulada \emph{EFI}. En ella se almacenan
        directamente los kernels de los SO o arrancadores (como GRUB)
        y sus archivos de configuración. UEFI conoce el formato FAT y
        carga el SO en base a esos archivos.
    \item Define en qué áreas de memoria el kernel encontrará
        información sobre la configuración del hardware.
    \item Cada kernel (o boot loader) es un archivo con extensión
        \texttt{.efi} compilado y enlazado para soportar el estándar.
    \item \textbf{Secure boot}: soporta firmas de kernels para
        verificar su integridad. El firmware permite almacenar en la
        EPROM certificados con las claves públicas de las empresas u
        organizaciones que los desarrollan.
    \item Soporta booteo remoto y tablas de partición de tamaño
        arbitrario basadas en el estándar \textbf{GPT} (\emph{Guided
        Partition Table}).
\end{itemize}

% =====================================================================
\section{Sistemas de archivos especiales}
% =====================================================================

Algunos SO modernos como GNU/Linux soportan sistemas de archivos cuyo
contenido \emph{no necesariamente se almacena en disco}. Se conocen
como \textbf{pseudo filesystems} o \emph{especiales}.

\subsection{procfs}

En GNU/Linux se usa \textbf{procfs}, comúnmente montado en
\verb|/proc|. Brinda información sobre datos del kernel y de cada
proceso en ejecución, en forma de un árbol de directorios y archivos.
Los usuarios pueden acceder y modificar datos (variables) internos del
kernel como si fueran archivos:

\begin{lstlisting}[style=shellstyle]
cat /proc/cpuinfo    # detalles de las CPUs
cat /proc/meminfo    # uso de la memoria
\end{lstlisting}

Salida típica de \verb|cat /proc/meminfo|:

\begin{lstlisting}[style=textstyle]
MemTotal:       994548 kB
MemFree:         65228 kB
MemAvailable:   263724 kB
Buffers:         21396 kB
Cached:         304440 kB
SwapCached:      25260 kB
Active:         267424 kB
Inactive:       503720 kB
Active(anon):   110956 kB
Inactive(anon): 351176 kB
...
\end{lstlisting}

Es posible modificar parámetros del kernel mediante:

\begin{lstlisting}[style=shellstyle]
echo <value> > /proc/sys/<class>/<parameter>
\end{lstlisting}

\subsection{sysfs}

\textbf{sysfs} está diseñado específicamente para contener archivos
especiales que representan dispositivos. Comúnmente se monta en
\verb|/dev|. Simplifica la manipulación de parámetros del kernel y de
dispositivos/drivers. En sistemas tipo BSD-UNIX o Linux se logra lo
mismo con el comando (y syscall) \texttt{sysctl}.

\subsection{ramfs e initramfs}

Otros filesystems se implementan en memoria principal, como
\textbf{ramfs}. Se usan durante el \emph{booting} como sistema de
archivos raíz inicial: el boot loader carga en memoria una imagen del
filesystem y se la presenta al kernel como un \textbf{initramfs}. Estos
filesystems contienen un conjunto reducido de programas de inicio y de
montaje de otros sistemas de archivos en disco.

\paragraph{busybox.}
En Linux, el \emph{initramfs} contiene comúnmente un FS con las
utilidades del paquete \textbf{busybox}, que contiene versiones
mínimas de las utilidades más comunes de UNIX en un único ejecutable.
Busybox también se usa en \emph{embedded systems}, ya que permite
generar un FS muy pequeño con algunos archivos de configuración en
\texttt{/etc}, dispositivos en \verb|/dev/| y un kernel (generalmente
Linux).

% =====================================================================
\section{Cierre}
% =====================================================================

El subsistema de archivos es la capa que organiza la persistencia y le
da forma humana al almacenamiento masivo. Las ideas más importantes:

\begin{itemize}[leftmargin=*]
    \item Un disco se ve a través de \textbf{volúmenes}; cada uno
        contiene un \emph{filesystem} con su propio formato (FAT,
        ext, NTFS, \ldots).
    \item Cada archivo es un par
        $(\textit{attributes},\;\textit{data})$ y se organiza dentro
        de un \textbf{árbol de directorios}. La unidad de
        identificación es el \emph{path}.
    \item El \textbf{VFS} provee un \emph{common file model} basado en
        \texttt{sb}, \texttt{file}, \texttt{inode}, \texttt{dir} y
        \texttt{fops}. Cada FS concreto mapea su formato a este
        modelo y se enchufa al kernel implementando una API definida.
    \item Dos formatos paradigmáticos: \textbf{FAT} (listas enlazadas
        de clusters indexadas en la tabla FAT) e \textbf{inodes} (cada
        archivo es un nodo índice con metadatos + tabla de bloques;
        para archivos grandes se usan inodes indirectos
        simple/doble/triple).
    \item El \textbf{montaje} cuelga otro FS en algún directorio del
        árbol raíz. La tabla de montaje del VFS le permite enrutar
        cada operación al filesystem correcto.
    \item Bajo el principio Unix de ``todo es un archivo'', los
        dispositivos viven como \textbf{archivos especiales} en
        \verb|/dev|, con un par (major, minor): major elige el
        driver, minor elige el dispositivo concreto. El descriptor de
        archivo del proceso (\verb|ofiles|) apunta a una
        \texttt{struct file} con cursor compartido.
    \item Los \textbf{dispositivos de caracteres} no usan caché de
        bloques: la E/S va directa al driver. Ejemplo paradigmático:
        la consola de xv6, abierta tres veces por \texttt{init} y
        heredada vía \texttt{fork}.
    \item El \textbf{filesystem de xv6} está organizado en capas
        (file descriptor / pathname / directory / inode / logging /
        buffer cache / disk driver). Es \emph{transaccional}: el log
        garantiza que un crash deje el FS en estado consistente.
    \item El \textbf{booting} arranca con el firmware (BIOS) que carga
        un boot loader (eventualmente multiboot, como GRUB); el
        estándar moderno UEFI usa una partición FAT con archivos
        \texttt{.efi} y soporta \emph{secure boot} y GPT.
    \item Los \textbf{filesystems especiales} (procfs, sysfs, ramfs,
        initramfs con busybox) exponen información o servicios del
        kernel como si fueran archivos: el mismo principio uniforme
        que con los dispositivos.
\end{itemize}

\begin{idea}
La gran abstracción del sistema de archivos es presentar a procesos y
usuarios un \emph{único árbol uniforme} en el que conviven discos,
particiones, dispositivos, parámetros del kernel y procesos en
ejecución. Esa uniformidad es la base que permite construir, encima,
un shell expresivo y herramientas Unix simples y componibles.
\end{idea}

\end{document}
